\documentclass[conference]{IEEEtran}
\usepackage[utf8]{inputenc}
\usepackage{hyperref}
\title{Effects of blue light exposure on sleep quality in young adults}
\author{\IEEEauthorblockN{ResearchGenie}\IEEEauthorblockA{AI-generated research draft}}
\begin{document}
\maketitle
\textit{\textbf{AI-generated research draft for human review.} Not a verified or published research paper. Every claim must be checked against its cited source before any external use.}\par

\section{Introduction}
\textit{Introduction generation did not satisfy evidence constraints.}

\section{Background}
The biological mechanism linking artificial light to sleep disruption centers on blue wavelengths activating melanopsin-containing retinal ganglion cells, which subsequently influence circadian rhythms and sleep onset [1]. This pathway is particularly relevant for young adults who frequently engage with electronic devices at night. While the core circadian system matures early in childhood, as noted regarding toddlers and pre-schoolers [2], the behavioral drivers of sleep disturbance shift in adolescence and adulthood. For instance, female university students are hypothesized to exhibit poorer sleep quality than males because they demonstrate a higher propensity for prolonged smartphone usage at bedtime [3]. This specific behavior leverages the blue light mechanism to suppress melatonin production, thereby disrupting circadian timing. Although this hypothesis addresses young adults, it draws upon broader models of how external light schedules induce disruptions in developing circadian systems [2]. Additionally, distinct environmental stressors like Ramadan intermittent fasting represent unique variants of metabolic restriction that align with specific cultural calendars, though they operate through different physiological pathways than light exposure [4].

\section{Literature Review}
The relationship between evening light exposure and sleep quality in young adults is robustly supported by empirical evidence, yet the magnitude of this effect varies significantly based on individual susceptibility and specific behavioral contexts. A systematic review by Silvani et al. (2022) indicates that approximately one-fifth of studies report blue light exposure decreases sleep quality, while roughly half found it reduced sleep efficacy, and slightly less than half noted increased sleep latency [5]. These findings suggest a consistent negative association across the literature, though the heterogeneity in study designs and populations complicates precise quantification. Interventions targeting this pathway demonstrate that modifying device usage can yield substantial improvements; for instance, a five-day evening restriction on smartphones among university students with low nomophobia resulted in large to very large gains in total sleep time (+45 minutes), sleep efficiency (+12\%), and reduced wake after sleep onset [6]. Crucially, these sleep improvements translated into enhanced cognitive recovery, evidenced by significant increases in morning reaction time and lower-limb reaction speed comparable to short-term training gains [6], alongside marked reductions in stress and anxiety levels among participants less fearful of smartphone separation [6]. Conversely, individuals with high nomophobia exhibited no meaningful sleep improvements during the same intervention and experienced acute elevations in stress, highlighting that psychological dependence on devices may modulate the physiological response to light restriction [6].

The mechanism linking light exposure to sleep disruption extends beyond simple device usage to include broader environmental light schedules. Modeling studies indicate that earlier sleep timing parameters render individuals inherently more susceptible to nighttime waking induced by external light, and even a 30-minute delay in evening lights-off time can trigger wake episodes [2]. Furthermore, low daytime light levels combined with irregular schedules, particularly weekend delays in morning and evening light exposure, markedly worsen sleep consolidation and induce carry-over effects that shift sleep timing throughout the week [2]. While reducing blue light emission from mobile phone screens has been shown to significantly improve subjective sleep quality scores in student populations [1], other large-scale analyses of university students have found no significant relationship between smartphone usage metrics, including screen time and filter usage, and sleep quality outcomes [3]. This discrepancy suggests that factors such as age, body mass index, obstructive sleep apnea, hormonal cycles, and sleep environment often outweigh the specific effects of blue light exposure in determining overall sleep quality [3]. Consequently, while blue light is a potent disruptor, its impact is not uniform across all young adults and is heavily confounded by comorbidities and lifestyle variables distinct from light exposure itself.

In contexts involving intermittent fasting, such as Ramadan, the interplay between circadian disruption and sleep quality becomes even more pronounced. Studies of young female athletes reveal that Ramadan intermittent fasting (RIF) leads to a gradual decline in sleep quality and duration, reaching its nadir during the fourth week of the month [4]. This deterioration is associated with adverse impacts on neuromuscular and cognitive abilities, including significant declines in performance observed during the final weeks of fasting [4]. Although dietary intake and body composition remained stable in these cohorts, the cognitive deficits persisted both during and after the fasting period, indicating that sleep disruption may be a primary driver of functional decline rather than caloric restriction alone [4]. Similar patterns emerge in elderly populations, where diurnal intermittent fasting impairs sleep quality and increases daytime sleepiness, with sedentary individuals experiencing more severe declines in nocturnal sleep duration compared to physically active counterparts [7]. These findings collectively underscore that while blue light exposure is a critical modifiable factor for sleep health in young adults, its effects are mediated by individual psychological traits, environmental schedules, and concurrent physiological stressors like fasting.

\section{Discussion}
\textit{This section audit could not be resolved; unsupported content was omitted.}

\section{Research Gap}
\subsection{Direct vs. Indirect Mechanisms of Blue Light-Induced Sleep Disruption}
The corpus presents a contradiction regarding the primary pathway by which blue light affects sleep. Paper [4] and Paper [5] support the direct biological mechanism where blue light suppresses melatonin, leading to reduced sleep quality. However, Paper [3] explicitly refutes this in its specific context of young adults, finding no relationship between smartphone usage (and implied blue light) and sleep quality. This discrepancy suggests that current literature fails to isolate the physiological effect of blue light from confounding behavioral variables.

\textit{Evidence:} Paper [4] states that reducing blue light 'leads to improved subjective quality of sleep' based on the theory that it blocks melatonin. Conversely, Paper [3] found 'no relationship between smartphone usage and sleep quality' despite the hypothesis that blue light disrupts circadian rhythms. The corpus does not contain studies that statistically control for cognitive hyperarousal or screen addiction to isolate the pure spectral effect of blue light.

\subsection{Standardization of Sleep Metrics Across Objective and Subjective Measures}
There is a lack of consensus on how to define 'sleep quality' in blue light research. The corpus utilizes both subjective indices (PSQI) and objective actigraphy but rarely integrates them to validate findings.

\textit{Evidence:} Paper [4] relies solely on the PSQI, reporting significant improvement after intervention. Paper [7] uses wrist-worn actigraphy to show large improvements in total sleep time and efficiency. However, no study in the corpus correlates these two metrics to determine if subjective perception aligns with objective physiological data when blue light is manipulated.

\subsection{Impact of Concurrent Environmental and Psychological Confounders}
Current studies often treat blue light as an isolated variable, ignoring significant confounders that may explain contradictory results or mask true effects.

\textit{Evidence:} Paper [7] highlights 'nomophobia' (fear of smartphone separation) as a factor limiting the effectiveness of digital restriction interventions. Paper [2], while focused on Ramadan, explicitly lists environmental variables (noise, household size) and emotional status (depression) as potential confounders that should be controlled in future investigations. The corpus lacks studies that simultaneously measure blue light exposure and these specific psychosocial/environmental factors.

\subsection{Sleep Disruption in Non-Student Young Adult Populations}
The evidence base is heavily skewed toward university students, limiting generalizability to young adults with different circadian demands.

\textit{Evidence:} Papers [3], [4], and [7] exclusively study university students (aged 17-22). Paper [5], a systematic review, notes that few resources summarize effects on healthy humans broadly, yet the retrieved corpus offers no data on young adults in shift-work roles or those with irregular schedules who may be more susceptible to blue light disruption.

\section{Proposed Novelty and Contribution}
Research addressing the direct physiological mechanism of blue light versus behavioral confounders in young adults represents high novelty.

\textit{Caveats:} Novelty is assessed relative to the provided corpus; it does not guarantee novelty against the entire global literature. The claim relies on the contradiction found between Paper [3] and Paper [4].

\section{Limitations}
Limitations in the current literature stem primarily from methodological constraints and narrow population scopes that restrict generalizability. Several studies rely on subjective self-report measures rather than objective tools; for instance, Boujelbane et al. (2022) utilized questionnaires like the PSQI instead of actigraphy or polysomnography to assess sleep quality [7], while Silvani et al. (2022) noted that most included studies exposed participants to blue light during non-strenuous tasks, leaving physical stress and hormonal influences unexplored [5]. Furthermore, many investigations lack control groups or fail to account for environmental variables; Bougrine et al. (2024) omitted a control group due to ethical considerations regarding fasting in Muslim countries and did not control for the timing or quantity of pre-dawn meals [4], whereas Chaudhry et al. (2024) lacked geographic specification and environmental controls, limiting broader applicability [3]. Population homogeneity further constrains findings; Alaya et al. (2026) restricted their sample to physically active students over a five-day period, precluding extrapolation to older adults or those with clinical sleep disorders [6], and Bougrine et al. (2024) focused exclusively on female athletes, ignoring potential sex-specific hormonal fluctuations [4]. These methodological gaps highlight the need for future research to expand recruitment to community-based settings, elucidate molecular mechanisms of the HPO axis under sleep deprivation, and explore therapeutic interventions such as melatonin supplementation or light therapy to reverse fertility impairment [8].

\section{Conclusion}
\textit{No evidence-backed conclusion can be drawn.}

\begin{thebibliography}{99}
\bibitem{ref1} [1] P. Randjelović, N. Stojanović, and I. Ilić, "The effect of reducing blue light from smartphone screen on subjective quality of sleep among students," \textit{Chronobiology international}, 2023, doi: 10.1080/07420528.2023.2173606.
\bibitem{ref2} [2] T. Yao and V. Booth, "External light schedules can induce nighttime sleep disruptions in a Homeostat-Circadian-Light Model for sleep in young children," \textit{arXiv (preprint)}, 2025.
\bibitem{ref3} [3] H. Chaudhry, H. Patel, and S. T. Avadhootha, "The Relationship Between Smartphone Usage and Sleep Quality Amongst University Students," \textit{arXiv (preprint)}, 2024.
\bibitem{ref4} [4] H. Bougrine, N. Chalghaf, and C. Azaiez, "The impact of intermittent fasting during Ramadan on psychomotor and cognitive skills in adolescent athletes," \textit{Frontiers in Sports and Active Living}, 2024, doi: 10.3389/fspor.2024.1362066.
\bibitem{ref5} [5] M. I. Silvani, R. Werder, and C. Perret, "The influence of blue light on sleep, performance and wellbeing in young adults: A systematic review," \textit{Frontiers in physiology}, 2022, doi: 10.3389/fphys.2022.943108.
\bibitem{ref6} [6] W. B. Alaya, M. A. Souissi, and I. Dergaa, "Effects of a five-day evening smartphone restriction on sleep, cognitive, and physical performance in university students: A single-arm repeated measures study stratified by nomophobia," \textit{Medicine}, 2026, doi: 10.1097/md.0000000000049698.
\bibitem{ref7} [7] M. A. Boujelbane, K. Trabelsi, and H. Jahrami, "Time-restricted feeding and cognitive function in sedentary and physically active elderly individuals: Ramadan diurnal intermittent fasting as a model," \textit{Frontiers in Nutrition}, 2022, doi: 10.3389/fnut.2022.1041216.
\bibitem{ref8} [8] N. J. Su, J. Huang, and H. Li, "Association between self-reported poor sleep status and ovulation rate in Chinese women with fertility intentions: a multicentre prospective cohort study protocol," \textit{BMJ open}, 2026, doi: 10.1136/bmjopen-2025-113145.
\end{thebibliography}

\end{document}
